\section{Model Komunikasi Aman dan Ancaman}

\subsection{Model Komunikasi Kriptografi}

Model standar melibatkan:
\begin{itemize}
  \item \textbf{Alice (Pengirim)}: Mengirim plaintext, melakukan enkripsi
  \item \textbf{Bob (Penerima)}: Menerima ciphertext, melakukan dekripsi
  \item \textbf{Eve (Penyerang)}: Mencoba mencuri dengar (eavesdropping) atau memodifikasi pesan
  \item \textbf{Kanal}: Saluran komunikasi (internet, udara, jaringan) yang mungkin tidak aman
\end{itemize}

Tujuan kriptografi adalah memastikan komunikasi aman meskipun kanal tidak terpercaya. Model komunikasi ini mencerminkan skenario dunia nyata di mana informasi sensitif harus ditransmisikan melalui jaringan publik yang rentan terhadap intervensi pihak ketiga. Dalam implementasi praktis, Alice dan Bob dapat merepresentasikan berbagai entitas seperti pengguna dan server, dua sistem dalam jaringan korporat, atau bahkan proses yang berkomunikasi dalam satu sistem komputer.

\subsection{Ancaman Keamanan}

\begin{table}[htbp]
\centering
\small
\begin{tabular}{|l|p{6cm}|}
\hline
\textbf{Ancaman} & \textbf{Deskripsi} \\
\hline
Interception & Eve mencuri dengar pesan tanpa diketahui \\
\hline
Modification & Eve mengubah isi pesan dalam perjalanan \\
\hline
Fabrication & Eve membuat pesan palsu seolah dari Alice \\
\hline
Replay & Eve merekam pesan dan mengirim ulang di waktu lain \\
\hline
Masquerade & Eve berpura-pura sebagai pihak lain \\
\hline
\end{tabular}
\caption{Jenis Ancaman Keamanan}
\end{table}

Ancaman keamanan diklasifikasikan berdasarkan dampaknya terhadap layanan keamanan informasi \cite{rfc4949}. Ancaman terhadap \textbf{confidentiality} mencakup serangan langsung untuk mendapatkan akses ke sistem atau teknik man-in-the-middle di mana penyerang memposisikan diri dalam aliran informasi untuk mencegat data. Ancaman terhadap \textbf{integrity} berupa modifikasi tidak sah terhadap data atau log sistem untuk menyembunyikan serangan.

{\sloppy Untuk mengatasi ancaman tersebut, kriptografi menyediakan mekanisme: \textbf{confidentiality} (kerahasiaan), \textbf{integrity} (keutuhan), \textbf{availability} (ketersediaan), dan \textbf{non-repudiation} (ketidakbantahan).\par} Ketiga layanan pertama dikenal sebagai CIA triad; non-repudiation melengkapi model untuk transaksi yang memerlukan bukti hukum. Implementasi yang tepat memastikan komunikasi digital tetap aman dan terpercaya.
